\newcommand{\RESULTSKEW}{%
Plain federated averaging stayed close to where it started. At the default $\alpha=0.5$ FedAvg reaches 0.800 on two seeds, the same as v0 (0.800 on three seeds), and well below centralised training on all data (0.850). The pairs per seed show why there is no gain to report: on seed 42 federation improved on v0 by 1.7 points (0.788 to 0.805), on seed 43 it lost 0.3 (0.798 to 0.795). With two seeds that is no detectable gain. With FedAvg, the 60\% of the training photos held by the farms added nothing that the 40\% central pool had not already given the model, under these settings (30 rounds, one local epoch, local learning rate $10^{-5}$ chosen in a short pilot of three values).

Label skew has a mild effect. From almost identical clients ($\alpha=100$) to strongly skewed ones ($\alpha=0.1$) the means fall from 0.794 to 0.779, but the $\alpha=100$ cell is not better than $\alpha=0.5$ and the two seeds differ from each other by as much as the cells do, so only the most skewed setting stands out, and only just. FedProx, which pulls each local model back towards the global one, does better. At $\alpha=0.5$ it reaches 0.811 on both seeds (0.812 and 0.810), above FedAvg on each seed (by 0.7 and 1.5 points) and above v0 on each seed (by 2.4 and 1.2 points). It is the only federated setting here that beats the central starting model on every seed it was run on. This rests on two seeds, so it is a consistent sign and not a firm ranking. At $\alpha=0.1$ FedProx was run on one seed only (0.801 against 0.797 for FedAvg on that seed), which cannot be separated from noise.}

\newcommand{\RESULTSCOPE}{%
Training and sending only the last blocks (430\,k parameters, 1.72 MB) gives 0.803 and the classifier head alone (20\,k parameters, 80 KB) gives 0.796, against 0.800 for the whole network (8.84 MB). The three differ by at most 0.007, about the seed-to-seed spread (0.003 to 0.005), so the cheapest scope gives up nothing measurable here. Over 30 rounds a client would upload about 265 MB for the full network and 2.4 MB for the head. For farms on slow rural connections this is the result that matters most: it makes the update small enough to send, at no cost in accuracy that these runs can detect.

The central stage is not optional. Starting from the ImageNet backbone with an empty head, the full network reaches only 0.287 after 30 rounds on both seeds (validation macro-F1 was still rising slowly, 0.30 at the last round). Federated learning with 10 clients, one local epoch and this learning rate does not rebuild what central pre-training provides in a few thousand images.}

\newcommand{\RESULTCLASS}{%
The tuber classes with enough photos are reliable (F1 at or above 0.93 for Blackspot Bruising, Brown Rot, Dry Rot, Healthy and Soft Rot). The weak spot is the leaf classes that look alike: Fungi, Pest and Phytophthora score 0.62 to 0.67, and the confusion matrix shows why. Pest is read as Fungi in 25 of 89 test photos, and Fungi as Phytophthora or Pest in 17 and 16 of 109. Two rare tuber classes, Blackleg and Common Scab, score 0.59 and 0.55 but have 7 test photos each, so those two scores can move by several points with one photo. This is why the service reports a tier and not just a label: 11 of the 16 held-out sample photos were ``very likely'' and correct, and the 4 wrong ones were all ``possible''.}

\newcommand{\CONCLUSION}{%
The deliverable is a working end-to-end system: a scripted pipeline from raw photos to a calibrated model, a containerised API that answers with a class, a confidence tier and a heat-map, drift monitoring on the inputs, tests, and instructions that let a reviewer check it in minutes (a mini-dataset smoke run, a 25-minute real federated run from the shipped starting model) or rerun the whole grid.

The experiments give three findings. First, centralised training on all the photos (0.850) is better than anything federated here: FedAvg gives about 0.80 and FedProx at $\alpha=0.5$ gives 0.811, so federation cost four to five points of macro-F1 and recovered at most one to two points of the five-point gap between the central starting model and centralised training. Second, label skew matters less than the choice of starting point; without the central stage the full network fails (0.287). Third, the cheapest update, the classifier head at 80 KB, is as good as the whole network at 8.84 MB, so communication cost does not have to be paid for accuracy.

What this does not show is how the system behaves on photos from real Kenyan farms, since the clients are simulated from two public datasets and the test set comes from the same sources. The drift monitor is the part of the system built for that gap.}
